% 大熊座（综述）
% license CCBYSA3
% type Wiki

本文根据 CC-BY-SA 协议转载翻译自维基百科\href{https://en.wikipedia.org/wiki/Ursa_Major}{相关文章}。

大熊座（Ursa Major），亦称“巨熊座”，是位于北天的一座星座，其相关神话很可能可以追溯至史前时代。其拉丁名称意为“较大的（或更大的）熊”，用以指代并区别于附近的小熊座（Ursa Minor）。$^{[2]}$ 在古代，大熊座是公元2世纪托勒密所列出的最初48个星座之一，其编制基础源自更早的希腊、埃及、巴比伦与亚述天文学家的研究成果。$^{[3]}$在现代被正式划定的88个星座中，大熊座按面积排名第三。

大熊座最为人熟知的是由其七颗主星构成的星群，这一星群在不同文化中被称为“大斗”“马车”“查理的车”或“犁”等名称。尤其值得注意的是，“大斗”的恒星排列形态在视觉上与“小斗”相呼应。其中的两颗恒星——天枢与天璇（分别为大熊座$\alpha$星与$\beta$星，Dubhe 与 Merak）——可作为导航指示星，用来指向当前的北极星，即位于小熊座中的北极星（Polaris）。

大熊座及其所包含或交叠的星群在世界诸多文化中具有重要意义，常被视为“北方”的象征。阿拉斯加州旗帜上对大熊座的描绘，便是这一象征在现代的一个实例。

在北半球的大多数地区，大熊座全年可见，并在中北纬度地区呈现为拱极星座。从南半球的温带纬度看，其主要星群不可见，但星座的南部区域仍可被观测到。
\subsection{特征}  
大熊座覆盖约1279.66平方度，约占全天空的3.10\%，是第三大星座。$^{[4]}$ 1930年，欧仁·德尔波特为其划定了国际天文学联合会（IAU）认可的正式星座边界，将其定义为一个由28条边组成的不规则多边形。在赤道坐标系中，该星座的赤经范围介于$08^{\mathrm h}08.3^{\mathrm m}$至$14^{\mathrm h}29.0^{\mathrm m}$之间，赤纬范围介于$+28.30^\circ$至$+73.14^\circ$ 之间。$^{[5]}$大熊座与八个星座相邻：北方与东北方为天龙座，东方为牧夫座，东与东南为猎犬座，东南为后发座，南方为狮子座与小狮座，西南为天猫座，西北为鹿豹座。国际天文学联合会于1922年采用其三字母缩写“UMa”。$^{[5]}$
\subsection{特征}  
\subsubsection{星群}
\begin{figure}[ht]
\centering
\includegraphics[width=6cm]{./figures/f2d88b83474133e7.png}
\caption{大熊座与北极星，以及北斗七星中明亮恒星的名称} \label{fig_UrMa_1}
\end{figure}
\begin{figure}[ht]
\centering
\includegraphics[width=6cm]{./figures/3eb25f5a42cf8f25.png}
\caption{肉眼所见的大熊座星座} \label{fig_UrMa_2}
\end{figure}
大熊座中七颗明亮恒星的轮廓构成了一个星群，在美国和加拿大被称为“北斗”（Big Dipper），而在英国则被称为“犁”（Plough）$^{[6]}$，或在历史上称为“查理之车”（Charles' Wain）$^{[7]}$。这七颗恒星中有六颗的视星等为二等或更亮，使其成为夜空中最为人熟知的恒星图案之一。$^{[8][9]}$ 正如这些常见名称所暗示的那样，其形状常被认为类似于一把勺子、一种农用犁具，或一辆马车。在大熊座的语境中，它们通常被绘制为巨熊的后躯与尾巴。从“勺斗”部分开始，沿顺时针方向（即天空中的东方方向）延伸至柄部，这些恒星依次为：
\begin{itemize}
\item 天枢（Dubhe，“熊”），视星等为1.79，是全天第35亮的恒星，也是大熊座中第二亮的恒星。  
\item 天璇（Merak，“熊的腰部”），视星等为2.37。  
\item 天玑（Phecda，“大腿”），视星等为2.44。  
\item 天权（Megrez），其名称意为“尾巴的根部”，指其位于熊的躯体与尾巴的交汇处（或北斗中斗身与柄部的连接点）。  
\item 玉衡（Alioth），其名称并非指熊，而是意为“黑马”，这是由原名讹变并误用于与其名称相近的伴星开阳（Alcor）所致，后者是肉眼可见的双星伴侣，围绕着开阳主星运行。$^{[10]}$ 玉衡是大熊座中最亮的恒星，也是全天第33亮的恒星，视星等为1.76。它同时也是化学丰度异常的 Ap 型恒星中最明亮的一颗，这类磁性恒星的某些化学元素会出现异常贫化或富集，并且随着恒星自转而表现出变化。$^{[10]}$
\item 开阳（Mizar，大熊座 $\zeta$ 星），位于北斗柄端向内数第二颗，是大熊座中第四亮的恒星。“Mizar”一名意为“腰带”。它与其视伴星辅星（Alcor，大熊座80号星）共同构成著名的双星系统，在阿拉伯传统中被称为“马与骑士”。  
\item 摇光（Alkaid，大熊座 $\eta$ 星），位于熊尾的末端。其视星等为1.85，是大熊座中第三亮的恒星。$^{[11][12]}$
\end{itemize}
除天枢与摇光之外，北斗七星中的其他恒星都具有指向人马座方向的共同自行运动。人们还识别出少数具有相同运动特征的其他恒星，它们共同被称为“大熊座移动星群”（Ursa Major Moving Group）。
\begin{figure}[ht]
\centering
\includegraphics[width=6cm]{./figures/58bc21206ce95031.png}
\caption{大熊座与小熊座相对于北极星的位置关系} \label{fig_UrMa_3}
\end{figure}